% ============================================================================
%  Threshold Selection and Sensitivity
%  Drop-in section for Chapter 3. Place after \label{sec:boundaries} and
%  before the chapter summary.
%
%  Sources for the figures in this section:
%   - Table \ref{tab:l3_sweep}: l3_threshold_sensitivity.py, 5 seeds,
%     5,005 attacks + 1,665 generated benign + 333 BIPIA benign.
%   - L4 figures: threshold_sensitivity_analysis.py and diagnose_fpr.py.
%   - Held-out FPR: build_fresh_holdout_fpr.py, seed 1, n = 333.
% ============================================================================

\section{Threshold Selection and Sensitivity}
\label{sec:thresholds}

Four numeric parameters govern how SecureRAG behaves. This section sets out what each one controls and the evidence behind the value chosen for it. Reporting them together makes the framework reproducible and allows the results in Chapter~\ref{ch:evaluation} to be interpreted against the settings that produced them.

\begin{table}[h]
\centering
\caption{The framework's numeric parameters.}
\label{tab:thresholds}
\begin{tabular}{llp{6.4cm}}
\toprule
\textbf{Parameter} & \textbf{Value} & \textbf{Basis for the value} \\
\midrule
\texttt{ANOMALY\_THRESHOLD} (L3) & $15.0$ & Sweep across attacks and both benign sets (Section~\ref{subsec:l3_threshold}) \\
\texttt{SEMANTIC\_THRESHOLD} (L4) & $0.18$ & Sweep bounded by measured internal similarity scores (Section~\ref{subsec:l4_threshold}) \\
L0 risk bands & $1.8\times$, $1.0\times$ & Multiples of the anomaly threshold (Section~\ref{subsec:l0_threshold}) \\
Retrieval relevance floor & $0.15$ & Relevance filter with a fallback; held fixed throughout (Section~\ref{subsec:retrieval_threshold}) \\
\bottomrule
\end{tabular}
\end{table}

\subsection{The anomaly threshold}
\label{subsec:l3_threshold}

\texttt{ANOMALY\_THRESHOLD} sets the scale for L3's blocking decision. L3 blocks a query when its anomaly score exceeds twice the effective threshold. The effective threshold is the nominal value for LOW and MEDIUM queries, and $0.7$ times that value for queries L0 has classified HIGH. The operative bar is therefore $30.0$ in the ordinary case and $21.0$ for high-risk queries.

Layers L0 through L3 read only the query text. They perform no retrieval and invoke no language model, so their decisions can be reproduced from the generators alone at negligible cost. This makes a full sweep practical, and Table~\ref{tab:l3_sweep} reports one across five seeds. Each candidate value is evaluated against the attack set and against both benign sets, so the effect on detection and on false positives can be read from the same row.

\begin{table}[h]
\centering
\caption{Detection and false-positive rate as a function of \texttt{ANOMALY\_THRESHOLD}. Five seeds; 5{,}005 attacks, 1{,}665 generated benign queries, 333 BIPIA benign documents. Detection is measured before L4.}
\label{tab:l3_sweep}
\begin{tabular}{rrrrr}
\toprule
\textbf{Threshold} & \textbf{Detection} & \textbf{L3 blocks} & \textbf{Internal FPR} & \textbf{External FPR} \\
\midrule
 5.0 & 93.33\,\% & 323 & 4.38\,\% & 91.29\,\% \\
 7.5 & 92.23\,\% & 268 & 0.00\,\% & 84.68\,\% \\
10.0 & 90.69\,\% & 191 & 0.00\,\% & 57.06\,\% \\
12.5 & 89.67\,\% & 140 & 0.00\,\% & 14.11\,\% \\
\textbf{15.0} & \textbf{89.37\,\%} & \textbf{125} & \textbf{0.00\,\%} & \textbf{1.80\,\%} \\
17.5 & 87.77\,\% &  45 & 0.00\,\% & 0.00\,\% \\
20.0 & 87.49\,\% &  31 & 0.00\,\% & 0.00\,\% \\
25.0 & 87.21\,\% &  17 & 0.00\,\% & 0.00\,\% \\
30.0 & 87.01\,\% &   7 & 0.00\,\% & 0.00\,\% \\
\bottomrule
\end{tabular}
\end{table}

The two sides of the sweep behave very differently. Detection declines gradually and almost linearly, losing 6.3 points across the entire range. The external false-positive rate does not: it falls from 91.29\,\% to 1.80\,\% between $5.0$ and $15.0$, then reaches zero immediately above. The shape of that curve is what selects the value. At $15.0$ the framework holds the external false-positive rate below 2\,\% while keeping the highest detection available at that level, and every step downward from there costs far more in false positives than it returns in detection.

The value $17.5$ is a reasonable alternative. It removes the remaining external false positives entirely and costs 1.6 points of detection. The results in this thesis were produced at $15.0$, and a deployment that weights false positives more heavily than missed detections would be well served by the higher value. Presenting both is the purpose of running the sweep.

The steep section of the curve between $12.5$ and $17.5$ has an identifiable cause. L3's first dimension scores length and sentence count, and long documents accumulate score quickly: a 90-sentence business email once scored $+108$ from sentence count alone. Measurement on the BIPIA benign set found that 211 of its 333 documents exceed 800 characters and contain no suspicious keyword at all. The dimension is now capped at a total contribution of $20.0$, which allows length to raise a query to MEDIUM risk and open the inexpensive L4 check while leaving it short of the hard-block bar. The six documents that remain blocked at $15.0$ are all dense data tables, and are analysed in Chapter~\ref{ch:evaluation}.

\subsection{The semantic threshold}
\label{subsec:l4_threshold}

\texttt{SEMANTIC\_THRESHOLD} is the cosine-similarity floor below which L4 treats a generated response as having left the semantic boundary of the knowledge base. The value was revised twice as measurement became available, and both revisions are worth recording because they show what the parameter is sensitive to.

The framework originally used $0.45$. Measured against the external benign set, that value produced a false-positive rate of $68.47\,\%$, rising to $96.67\,\%$ on email documents. Inspecting 29 of those cases showed similarities between $0.223$ and $0.380$: correct answers, grounded in document content the query itself supplied, that simply did not resemble the corpus. The value was also out of step with the retrieval stage, which admits a passage as relevant at a similarity above $0.15$. Requiring an answer to resemble the corpus more closely than a passage must in order to be considered relevant material sets the output bar higher than the input bar, which is not a defensible ordering. The threshold was lowered to $0.15$ to bring the two into line.

The shipped value of $0.18$ comes from a sweep over both sides. On the external BIPIA data the false-positive rate is flat at $1.80\,\%$ for every threshold from $0.15$ through $0.25$, which suggests that raising the bar to $0.25$ would be free. The internal data shows otherwise. Five of the 333 internal benign responses score between $0.186$ and $0.2331$, and they include the kind of query the L0 correction was made to protect, such as questions about overriding a method in a Python subclass or about \texttt{.gitignore} configuration. A threshold of $0.25$ would block those at L4, undoing at the output stage a correction made at the input stage. Setting the threshold at $0.18$ leaves the lowest of those scores, $0.186$, clear of the bar. It is the highest value that costs nothing on either side, and any increase beyond it becomes a trade between external detection and internal false positives rather than a free gain.

\subsection{The L0 risk bands}
\label{subsec:l0_threshold}

L0 classifies a query as HIGH when its anomaly score exceeds $1.8$ times the anomaly threshold, MEDIUM between $1.0$ and $1.8$ times, and LOW below that. Neither boundary is an independent parameter. Both are multiples of the single value examined in Section~\ref{subsec:l3_threshold}, so calibrating that value fixes both.

The sensor's fast path for high-risk queries calls the same regular expressions L2 uses rather than keeping a keyword list of its own. Sharing L2's patterns keeps the two components aligned and prevents ordinary technical vocabulary from being treated as an attack simply because it contains a sensitive word.

The design also limits what a misjudgement at this stage can cost. L0 does not block. It returns a classification that adjusts two later decisions, lowering L3's effective bar for HIGH-risk queries and opening L4's gate, and no query is refused on L0's verdict alone. A query the sensor rates too highly is examined more closely rather than rejected, which is why L0 introduces no threshold requiring separate calibration.

\subsection{The retrieval relevance floor}
\label{subsec:retrieval_threshold}

The retriever admits a passage as context when its cosine similarity to the query exceeds $0.15$. Passages that rank in the top five but fall below this value are dropped, so that weakly related material does not enter the model's context window and enlarge the surface an injected instruction can arrive through.

The threshold acts as a relevance filter rather than a hard gate. If no retrieved passage clears it, the two highest-scoring passages are supplied regardless, so a legitimate query is never left without context. Its influence on the reported results is indirect: it determines which passages the model receives, and affects blocking only through the similarity L4 subsequently measures on the response. The value is held fixed across every result reported in this thesis. Establishing its sensitivity would require re-running retrieval and generation at each candidate value, and is identified as future work in Chapter~\ref{ch:conclusion}.

\subsection{Held-out measurement}
\label{subsec:threshold_circularity}

The 333-document BIPIA benign set used in Section~\ref{subsec:l3_threshold} also served earlier in the project to set two parameters: L3's length cap and L4's semantic threshold. A false-positive rate measured on that set is therefore a training-set figure for those two parameters, and a held-out figure is needed alongside it.

The raw BIPIA pool holds approximately 1{,}100 documents, of which 333 were used. Drawing a fresh sample of 333 from the 767 that were never used gives a false-positive rate of $1.20\,\%$, or 4 documents. The draw is seeded, so further independent samples are available on demand. The $1.20\,\%$ figure is the held-out result; the $1.80\,\%$ reported elsewhere is measured on the tuning set, and the two should be read with that distinction in mind.
